\section{Elementos de la pantalla}
\label{sec:pantalla}
\begin{description}
\item[Barra lateral.] A la izquierda. Tiene la lista de páginas y, debajo, un resumen de
la red activa o de la cartera activa. Se oculta con la flecha \texttt{<<}.
\item[Deslizador.] Barra con un punto que se arrastra para elegir un número. El valor
actual aparece encima.
\item[Formulario.] Recuadro con varios controles y un botón azul (\boton{Construir},
\boton{Simular}, \boton{Simular cartera}). Los cambios no se aplican hasta presionar el
botón.
\item[Métrica.] Número grande con su nombre encima. El símbolo \textbf{?} muestra una
explicación al pasar el cursor. Debajo puede aparecer la diferencia frente a un caso de
referencia.
\item[Gráfica interactiva.] Muestra información al pasar el cursor. Se amplía
arrastrando un rectángulo; doble clic restablece la vista. Los botones de la esquina
permiten descargar la imagen. Un clic en la leyenda oculta o muestra una serie.
\item[Tabla.] Un clic en el nombre de una columna la ordena.
\item[Desplegable.] Línea con una flecha \texttt{>} que muestra contenido adicional.
\item[Indicador de ejecución.] Aparece arriba a la derecha mientras la aplicación calcula.
\end{description}

\subsection{Estado compartido y límites}
\begin{itemize}
\item La red construida en \emph{Construir red} y los parámetros elegidos en
\emph{Diagnóstico}, \emph{Ritmo de pagos} y \emph{Simulador de choques} se comparten entre
las páginas 1 a 7. La cartera de \emph{Cartera multicapa} se comparte con la página 9.
\item El estado dura mientras la pestaña del navegador esté abierta. Al recargarla, todo
vuelve a los valores iniciales.
\item Con más de 300 entidades se desactivan las simulaciones en vivo; el \emph{Mapa de
riesgo} sigue funcionando porque predice sin simular. El límite está en
\texttt{configs/default.yaml} (\texttt{app.max\_live\_entities}).
\end{itemize}

\subsection{Conceptos mínimos}
\label{sec:conceptos}
\begin{description}
\item[Entidad o nodo.] Participante de la red: un banco, una empresa o una persona.
\item[Arista.] Pago recurrente de una entidad a otra. Su peso es el pago medio por ciclo.
\item[Ciclo.] Unidad de tiempo de las páginas 1 a 7 (un día hábil). En las páginas 8 y 9
la unidad es el mes.
\item[Déficit.] Liquidez que le falta a una entidad.
\item[Colchón.] Déficit máximo que soporta antes de dejar de pagar.
\item[Impago y cascada.] Una entidad deja de pagar; si eso lleva a otras a dejar de pagar,
se produce una cascada.
\item[Fase y sincronía.] Momento del ciclo de pagos en que está una entidad, y qué tan
alineados están los ciclos de todas.
\item[PD.] Probabilidad de impago de un cliente en los próximos doce meses.
\item[Semilla.] Número que fija el azar: con la misma semilla se obtiene el mismo
resultado.
\end{description}
